%! The Normal Distribution and z-Scores — Unit 1, Lesson 1.7 — Homework
% EFFL, AP Statistics variant: this is the lesson's SCORED practice and the LAST component in
% the packet. These items were the in-class Check Your Understanding set; they are now the
% graded homework, started in the closing minutes of the period and finished at home. The
% AP Practice component that precedes it stays assigned-but-unscored.
% This is the LAST content lesson of Unit 1, so the closing spiralbox previews the UNIT TEST
% and Unit 2 rather than a next lesson.
% \answerspace{H}{} reserves exactly H in the blank; the key passes the answer as the second
% argument, so the two files paginate identically by construction. Keep the macro and every
% \boxguard byte-identical in the blank and the key.
% Multi-step solutions go in a \begin{work} block authored BYTE-IDENTICALLY here and in
% the _key: the blank reserves its exact height, the key prints it.
\documentclass[10pt]{article}
\usepackage{apstats-article}
\usepackage{apstats-boxes}
\newcommand{\answerspace}[2]{\par\nopagebreak\noindent\begin{minipage}[t][#1][t]{\linewidth}%
  \color{keyred}\bfseries #2\end{minipage}\par}

\begin{document}

\pageheader{Unit 1, Lesson 1.7}{Homework}
% NAMESTRIP: no \namedateperiod here — the name row lives on the cover only.

\vspace{0.15in}

% Authored BYTE-IDENTICALLY in homework/ and homework_key/ — this box is what keeps the
% blank and the key the same number of pages.
\boxguard[8]
\begin{remindbox}
\small
\textbf{This is your graded homework.} It is scored and due next class. You will start it in the
last minutes of the period, so ask about anything that stalls you before you leave, then finish
it on your own.
\end{remindbox}

\vspace{0.12in}

% ── The scored practice (formerly this lesson's Check Your Understanding) ────
\boxguard[14]
\begin{notesbox}{Standardized Scores and the Normal Model}
Show the arithmetic wherever a question asks for a number, and write every explanation
\emph{in the context of the problem} --- that is what earns credit on the AP exam, and it is the
standard here.
\begin{enumerate}[leftmargin=1.9em, itemsep=6pt]
  \item Label each summary a \textbf{parameter} or a \textbf{statistic}.
        \begin{enumerate}[label=\textbf{(\roman*)}, leftmargin=2.3em, itemsep=4pt]
          \item The mean weight of all 4,200 loaves a bakery baked last year is 799 g.
                \blank{3.0cm}
          \item A health inspector weighs 40 of those loaves and finds a mean weight of 803 g.
                \blank{3.0cm}
          \item The standard deviation of the heights of all 1,150 students enrolled at a school
                is 8.4 cm. \blank{3.0cm}
          \item Of the school's 400 statistics students, a teacher samples 25 and finds their
                median quiz score is 18 points. \blank{3.0cm}
        \end{enumerate}

  \item Resting heart rate for adult women is approximately normal with $\mu = 74$ beats per
        minute and $\sigma = 8$ beats per minute. Find the $z$-score of each value.
        \begin{work}
        z_{90} &= \frac{90 - 74}{8} = \frac{16}{8} = 2.0 \\[3pt]
        z_{74} &= \frac{74 - 74}{8} = \frac{0}{8} = 0 \\[3pt]
        z_{62} &= \frac{62 - 74}{8} = \frac{-12}{8} = -1.5
        \end{work}
        Interpret the \emph{negative} one in a sentence, in context.
        \answerspace{2.4cm}{}

  \item Ana scored \textbf{82} on a chemistry test where the class mean was 74 points with a
        standard deviation of 4 points. Ben scored \textbf{88} on a history test where the class
        mean was 79 points with a standard deviation of 6 points.
        \begin{work}
        z_{\text{Ana}} &= \frac{82 - 74}{4} = \frac{8}{4} = 2.0 \\[3pt]
        z_{\text{Ben}} &= \frac{88 - 79}{6} = \frac{9}{6} = 1.5
        \end{work}
        Ben's raw score is higher. Who did better \emph{relative to their own class}? Justify
        your answer, and say in one sentence why the raw scores alone cannot settle it.
        \answerspace{3.0cm}{}

  \item A school bus route's travel time is approximately normal with $\mu = 32$ minutes and
        $\sigma = 3$ minutes. The curve is marked at the mean and at one, two, and three standard
        deviations either side of it.

        \vspace{4pt}
        \begin{center}
        \begin{tikzpicture}[scale=0.8]
          \draw[navy, very thick] plot[smooth] coordinates
            {(-3.5,0.01) (-3,0.03) (-2.5,0.11) (-2,0.34) (-1.5,0.81) (-1,1.52) (-0.5,2.21)
             (0,2.50) (0.5,2.21) (1,1.52) (1.5,0.81) (2,0.34) (2.5,0.11) (3,0.03) (3.5,0.01)};
          \foreach \x/\h in {-3/0.03, -2/0.34, -1/1.52, 0/2.50, 1/1.52, 2/0.34, 3/0.03} {
            \draw[linegray, thin, dashed] (\x,0) -- (\x,\h);}
          \draw[->, thick] (-4.0,0) -- (4.0,0);
          \foreach \x/\lab in {-3/23, -2/26, -1/29, 0/32, 1/35, 2/38, 3/41} {
            \draw (\x,-0.08) -- (\x,0.08); \node[below, font=\scriptsize] at (\x,-0.1) {\lab};}
          \node[font=\scriptsize] at (0,-0.62) {Travel time (minutes)};
        \end{tikzpicture}
        \end{center}
        \vspace{2pt}

        \textbf{(i)} About what percent of trips take between 29 and 35 minutes? \blank{2.6cm}

        \vspace{4pt}
        \textbf{(ii)} About what percent take between 26 and 38 minutes? \blank{2.6cm}

        \vspace{4pt}
        \textbf{(iii)} About what percent take \emph{longer than} 38 minutes? \blank{2.6cm}
        \quad Explain how you got it.
        \answerspace{2.4cm}{}

  \item \textbf{AP-style multiple choice.} In Race X, finishing times are approximately normal
        with mean 21.0 minutes and standard deviation 1.5 minutes; runner P finished in
        18.0 minutes. In Race Y, times are approximately normal with mean 25.0 minutes and
        standard deviation 2.0 minutes; runner Q finished in 22.0 minutes. Which statement is
        \textbf{correct}? Circle one, then justify it in a sentence.
        \begin{enumerate}[label=\textbf{(\Alph*)}, leftmargin=2.2em, itemsep=2pt]
          \item Runner Q, because a standardized score of $-1.5$ is greater than $-2.0$.
          \item Runner P, because a standardized score of $-2.0$ puts P further below the mean of
                P's own race, and in a race a lower time is better.
          \item Runner P, because 18.0 minutes is the smaller raw finishing time.
          \item Runner Q, because Race Y had the larger standard deviation.
          \item Neither, because the two races had different means and cannot be compared.
        \end{enumerate}
        \answerspace{2.6cm}{}
\end{enumerate}
\end{notesbox}

\vspace{0.12in}

% The next-lesson preview closes the packet, so it lives here rather than in AP Practice.
% Lesson 1.7 is the LAST content lesson of Unit 1, so this looks ahead to the unit test and
% to Unit 2 — in plain, spoiler-free language.
\boxguard[12]
\begin{spiralbox}
\small
\textbf{Next: the Unit 1 test, then Unit 2 --- Exploring Two-Variable Data.} The test covers
everything in this packet's unit: telling categorical variables from quantitative ones and
choosing a display for each; describing a distribution's shape, centre, spread, and unusual
values \emph{in context}; mean, median, standard deviation, range, IQR, and which of those a
single extreme value can move; the five-number summary, boxplots, and comparing two groups; and
today's normal model, standardized scores, and the empirical rule.

\vspace{3pt}
Then everything changes shape. Every question so far has been about \emph{one} measurement at a
time. In Unit 2 we record \emph{two} measurements on the same individual --- a student's study
hours and that same student's test score --- and ask a new kind of question: when one of them
goes up, does the other tend to go up too? And if it does, how strongly, and can you use one to
predict the other?
\end{spiralbox}

\end{document}
